\documentclass[11pt]{article}

\usepackage[margin=1in]{geometry}
\usepackage{amsmath,amssymb,amsthm,mathtools}
\usepackage{enumitem}
\usepackage{microtype}
\usepackage[hidelinks]{hyperref}

\newtheorem{theorem}{Theorem}[section]
\newtheorem{proposition}[theorem]{Proposition}
\newtheorem{corollary}[theorem]{Corollary}
\newtheorem{lemma}[theorem]{Lemma}
\theoremstyle{definition}
\newtheorem{definition}[theorem]{Definition}
\newtheorem{example}[theorem]{Example}
\theoremstyle{remark}
\newtheorem{remark}[theorem]{Remark}

\newcommand{\R}{\mathbb{R}}
\newcommand{\simplex}{\Delta_{n-1}}
\newcommand{\supp}{\operatorname{supp}}
\newcommand{\diag}{\operatorname{diag}}

\title{Boundary Contact Order Governs Stationarity Preservation in Normalized Simplex Lifts}
\author{}
\date{}

\begin{document}

\maketitle

\begin{abstract}
Normalized nonnegative parametrizations convert optimization over a simplex into unconstrained smooth optimization, but their boundary singularities can hide the missing Karush--Kuhn--Tucker inequalities. We introduce the contact order of an inactive simplex coordinate and prove an exact finite-jet characterization. For a convex differentiable objective, a lifted point is critical through order $q$ precisely when its active partial derivatives agree and every inactive coordinate with contact order at most $q$ satisfies its missing KKT inequality. Consequently, preservation of simplex stationarity for arbitrary smooth objectives occurs at the sharp threshold given by the largest inactive contact order. Applied to deep Hadamard-product parametrizations, an inactive coordinate with $z$ zero factors has contact order $2z$. Ordinary second-order stationarity therefore sees exactly the coordinates with one zero factor, while coordinates with several zero factors require correspondingly higher jets. The same analysis produces exponentially many nonglobal second-order critical images for higher even-power lifts of a strongly convex quadratic. These results isolate a geometric quantity that determines which constrained optimality conditions a smooth lift can detect at finite order.
\end{abstract}

\section{Introduction}

Let
\[
  \simplex=\left\{x\in\R^n_+:\sum_{i=1}^n x_i=1\right\}.
\]
A common way to remove the simplex constraint is to optimize the composition
$F=f\circ\Phi$, where $f$ is defined near $\simplex$ and $\Phi$ is a
normalized nonnegative map. The squared lift
\[
  \Phi_i(u)=\frac{u_i^2}{\sum_k u_k^2}
\]
is the best-known example. At a boundary point its Jacobian loses the
directions that add a previously inactive coordinate. Second derivatives
recover those directions and expose the corresponding KKT inequalities.

This paper studies the general block-separable lift
\begin{equation}\label{eq:lift}
  \Phi_i(w)=\frac{a_i(w_i)}{A(w)},
  \qquad
  A(w)=\sum_{k=1}^n a_k(w_k),
\end{equation}
where the blocks $w_i$ are independent and every $a_i$ is smooth and
nonnegative. The central quantity is the first nonzero Taylor order of
$a_i$ at an inactive block. We call it the \emph{contact order}. It determines
the first derivative order at which the lifted objective can test the
missing KKT inequality for that coordinate.

The main contributions are as follows.
\begin{enumerate}[label=(\roman*)]
  \item For convex differentiable objectives, we give an exact
  characterization of finite-jet criticality in terms of active-gradient
  equality and the KKT inequalities whose contact orders are visible to the
  chosen jet.
  \item For arbitrary smooth objectives, we identify a sharp universal
  stationarity-preservation threshold: the jet order must reach the largest
  inactive contact order.
  \item For deep squared Hadamard products, we prove that contact order equals
  twice the number of zero factors. Thus second-order stationarity detects a
  missing coordinate exactly when it has one zero factor.
  \item For tied even-power lifts, we exhibit an exponential family of
  nonglobal second-order critical images for a strongly convex quadratic.
\end{enumerate}

The square-root simplex reparametrization and its second-order relation to
simplex stationarity are established in prior work
\cite{LiMcKenzieYin2023,LevinKileelBoumal2024}. General landscape-transfer
results explain when first- and second-order critical points survive smooth
parametrization \cite{LevinKileelBoumal2024}. Recent work also develops
Hadamard parametrizations for structured optimization
\cite{TangToh2024,OuyangEtAl2024,KolbEtAl2026}. The result here concerns higher
finite jets at boundary points and gives an exact order-by-order test for
normalized block-separable simplex lifts. Section~\ref{sec:related} states the
novelty boundary more explicitly.

\section{Setup and finite-jet criticality}\label{sec:setup}

For each $i$, let $E_i$ be a finite-dimensional real vector space and let
$U_i\subseteq E_i$ be open. Put $U=\prod_i U_i$. Fix an integer $Q\ge 1$.
Assume throughout that
\[
  a_i\in C^Q(U_i),\qquad a_i\ge 0,
\]
and consider the open set $U^+=\{w\in U:A(w)>0\}$. The map \eqref{eq:lift}
is $C^Q$ on $U^+$ and takes values in $\simplex$.

Fix $\bar w\in U^+$ and write
\[
  \bar x=\Phi(\bar w),\qquad
  S=\supp(\bar x)=\{i:a_i(\bar w_i)>0\},
  \qquad A_0=A(\bar w).
\]
For active blocks we impose the nondegeneracy condition
\begin{equation}\label{eq:active-nondegenerate}
  Da_i(\bar w_i)\ne 0 \qquad (i\in S).
\end{equation}
This ensures that independent block perturbations generate all tangent
directions within the active face.

\begin{definition}[Finite contact order]\label{def:contact}
For $j\notin S$, the function $a_j$ has \emph{contact order} $r_j$ at
$\bar w_j$ if $2\le r_j\le Q$,
\[
  D^k a_j(\bar w_j)=0\quad(1\le k<r_j),
  \qquad D^{r_j}a_j(\bar w_j)\ne 0.
\]
Equivalently, some $v_j\in E_j$ and some $c_j>0$ satisfy
\begin{equation}\label{eq:contact-path}
  a_j(\bar w_j+t v_j)=c_jt^{r_j}+o(t^{r_j}).
\end{equation}
We assume every inactive block under discussion has such a finite order.
\end{definition}

Nonnegativity forces each $r_j$ to be even. Indeed, if
$P(v)=D^{r_j}a_j(\bar w_j)[v^{r_j}]/r_j!$ is the first nonzero homogeneous
Taylor term, then $P(v)\ge0$ for all $v$. A nonzero odd homogeneous polynomial
cannot be nonnegative because $P(-v)=-P(v)$.

Our criticality definition records the Taylor sign along every curve. This
form is convenient because the first visible boundary direction can occur at
an order larger than two.

\begin{definition}[$q$-jet criticality]\label{def:jet}
Let $G$ be $C^q$ near $\bar w$. We call $\bar w$ \emph{$q$-jet critical} for
$G$ if, for every $C^q$ curve $\gamma:(-\epsilon,\epsilon)\to U^+$ with
$\gamma(0)=\bar w$, either
\[
  (G\circ\gamma)^{(k)}(0)=0\qquad(1\le k\le q),
\]
or the smallest $k\le q$ for which this derivative is nonzero is even and
positive.
\end{definition}

For $q=1$, Definition~\ref{def:jet} is ordinary first-order criticality. For
$q=2$, it is equivalent to zero gradient and positive-semidefinite Hessian.
The definition is necessary, not sufficient, for a local minimum when
$q>2$.

\begin{proposition}[Coordinate invariance]\label{prop:invariance}
Let $\psi_i:\widetilde U_i\to U_i$ be a $C^Q$ local diffeomorphism with
$\psi_i(\bar z_i)=\bar w_i$, and put $\psi=(\psi_1,\ldots,\psi_n)$. Then
$a_i\circ\psi_i$ has the same contact order at $\bar z_i$ as $a_i$ has at
$\bar w_i$. Moreover, $\bar w$ is $q$-jet critical for $F$ if and only if
$\bar z$ is $q$-jet critical for $F\circ\psi$.
\end{proposition}

\begin{proof}
Let $r$ be the first nonzero Taylor order of $a_i$ at $\bar w_i$. The first
nonzero homogeneous term of $a_i\circ\psi_i$ is
\[
  v\longmapsto
  \frac1{r!}D^r a_i(\bar w_i)
  [D\psi_i(\bar z_i)v]^r.
\]
The derivative $D\psi_i(\bar z_i)$ is invertible, so this polynomial is
nonzero and every lower term vanishes. Thus the contact order is unchanged.
The second assertion follows because a local $C^q$ diffeomorphism and its
inverse give a bijection between the germs of $C^q$ curves through the two
base points, without changing the scalar composition along corresponding
curves.
\end{proof}

Let $f$ be differentiable near $\simplex$ and put
\[
  g=\nabla f(\bar x),\qquad g_i=\partial_i f(\bar x).
\]
The simplex KKT stationarity condition at $\bar x$ is
\begin{equation}\label{eq:kkt}
  g_i=\lambda\quad(i\in S),
  \qquad
  g_j\ge\lambda\quad(j\notin S)
\end{equation}
for some $\lambda\in\R$.

\section{The contact-order theorem}\label{sec:main}

\begin{theorem}[Exact convex finite-jet characterization]\label{thm:convex}
Let $f$ be convex and $C^q$ on a neighborhood of $\simplex$, where
$1\le q\le Q$. Suppose \eqref{eq:active-nondegenerate} holds and every
$j\notin S$ has a finite contact order $r_j\le Q$. Then $\bar w$ is
$q$-jet critical for $F=f\circ\Phi$ if and only if there exists
$\lambda\in\R$ such that
\begin{align}
  g_i&=\lambda &&(i\in S),\label{eq:active-equality}\\
  g_j&\ge\lambda &&(j\notin S\text{ with }r_j\le q).
  \label{eq:visible-inequality}
\end{align}
No condition is imposed on an inactive index with $r_j>q$.
\end{theorem}

The theorem gives a complete filtration of the missing KKT inequalities. As
$q$ increases, an inactive coordinate enters the stationarity test exactly at
its contact order.

Convexity is needed only for the sufficiency direction of the exact
finite-jet characterization. The universal threshold for recovering
first-order constrained stationarity does not require convexity.

\begin{theorem}[Sharp universal stationarity threshold]\label{thm:threshold}
Let $f$ be $C^q$ near $\simplex$ and retain the hypotheses on the lift. If
$S\ne\{1,\ldots,n\}$ and
\[
  q\ge r_{\max}:=\max_{j\notin S}r_j,
\]
then $q$-jet criticality of $\bar w$ for $f\circ\Phi$ implies simplex KKT
stationarity of $\bar x$ for $f$.

The threshold is sharp over smooth objectives: if $q<r_j$ for some inactive
$j$, the linear objective $f(x)=-x_j$ makes $\bar w$ $q$-jet critical while
$\bar x$ violates simplex KKT stationarity. If $S=\{1,\ldots,n\}$, first-order
criticality already implies stationarity.
\end{theorem}

\section{Proofs}\label{sec:proofs}

We begin with the first-order calculation.

\begin{lemma}[Active-face gradient equality]\label{lem:first-order}
Under \eqref{eq:active-nondegenerate}, first-order criticality of $\bar w$ for
$F=f\circ\Phi$ is equivalent to
\[
  g_i=\lambda\qquad(i\in S)
\]
for some $\lambda$.
\end{lemma}

\begin{proof}
For any block direction $h=(h_1,\ldots,h_n)$, differentiation of
\eqref{eq:lift} gives
\begin{equation}\label{eq:first-variation}
  DF(\bar w)[h]
  =\frac{1}{A_0}\sum_{i=1}^n
    (g_i-g^\top\bar x)Da_i(\bar w_i)[h_i].
\end{equation}
For $j\notin S$, nonnegativity and $a_j(\bar w_j)=0$ imply
$Da_j(\bar w_j)=0$. Since the blocks are independent and
$Da_i(\bar w_i)\ne0$ for $i\in S$, the right-hand side vanishes for every
$h$ exactly when
\[
  g_i=g^\top\bar x\qquad(i\in S).
\]
Taking $\lambda=g^\top\bar x$ proves the claim. Conversely, this equality
makes every term in \eqref{eq:first-variation} vanish.
\end{proof}

\begin{lemma}[Boundary jet identity]\label{lem:boundary-jet}
Suppose active-gradient equality holds with value $\lambda$. For an inactive
$j$, choose $v_j$ as in \eqref{eq:contact-path} and perturb only the $j$th
block:
\[
  \gamma_j(t)=(\bar w_1,\ldots,\bar w_j+t v_j,\ldots,\bar w_n).
\]
Then
\begin{equation}\label{eq:jet-identity}
  (F\circ\gamma_j)^{(r_j)}(0)
  =\frac{D^{r_j}a_j(\bar w_j)[v_j^{r_j}]}{A_0}
   (g_j-\lambda).
\end{equation}
All lower derivatives vanish.
\end{lemma}

\begin{proof}
Write $h(t)=a_j(\bar w_j+t v_j)$. Along the path,
\[
  \Phi(\gamma_j(t))
  =\frac{A_0}{A_0+h(t)}\bar x
   +\frac{h(t)}{A_0+h(t)}e_j.
\]
Hence
\[
  \Phi(\gamma_j(t))-\bar x
  =\frac{h(t)}{A_0+h(t)}(e_j-\bar x)
  =\frac{c_j}{A_0}(e_j-\bar x)t^{r_j}+o(t^{r_j}).
\]
Differentiability of $f$ yields
\[
  F(\gamma_j(t))-F(\bar w)
  =\frac{c_j}{A_0}g^\top(e_j-\bar x)t^{r_j}
   +o(t^{r_j}).
\]
Active-gradient equality gives $g^\top\bar x=\lambda$, while
$r_j!c_j=D^{r_j}a_j(\bar w_j)[v_j^{r_j}]$. This proves
\eqref{eq:jet-identity} and the lower-order vanishing.
\end{proof}

We also need a curvewise estimate for coordinates whose contact order is
beyond the visible jet.

\begin{lemma}[Inactive-coordinate order]\label{lem:inactive-order}
If $a_j$ has contact order $r_j$ at $\bar w_j$ and
$\gamma_j(t)=\bar w_j+O(t)$, then
\[
  a_j(\gamma_j(t))=O(|t|^{r_j}),
  \qquad
  \Phi_j(\gamma(t))=O(|t|^{r_j}).
\]
\end{lemma}

\begin{proof}
Taylor's theorem and the vanishing of all derivatives below $r_j$ give
$a_j(\bar w_j+u)=O(\lVert u\rVert^{r_j})$. Every differentiable curve obeys
$\gamma_j(t)-\bar w_j=O(t)$. The denominator $A(\gamma(t))$ remains bounded
away from zero near $t=0$ because $A_0>0$, which proves the second estimate.
\end{proof}

\begin{proof}[Proof of Theorem~\ref{thm:convex}]
Assume first that $\bar w$ is $q$-jet critical. Since $q\ge1$,
Lemma~\ref{lem:first-order} gives active-gradient equality. For every inactive
$j$ with $r_j\le q$, Lemma~\ref{lem:boundary-jet} and the evenness of $r_j$
show that $g_j-\lambda$ cannot be negative. Thus
\eqref{eq:visible-inequality} holds.

Conversely, assume \eqref{eq:active-equality} and
\eqref{eq:visible-inequality}. Let $\gamma$ be any $C^q$ curve through
$\bar w$, write $x(t)=\Phi(\gamma(t))$, and use convexity at $\bar x$:
\begin{align}
  F(\gamma(t))-F(\bar w)
  &\ge g^\top(x(t)-\bar x)\\
  &=\sum_{j\notin S}(g_j-\lambda)x_j(t).
  \label{eq:convex-bound}
\end{align}
The terms with $r_j\le q$ are nonnegative. By
Lemma~\ref{lem:inactive-order}, each remaining term is
$O(|t|^{r_j})=O(|t|^{q+1})$. Therefore
\begin{equation}\label{eq:lower-order-bound}
  F(\gamma(t))-F(\bar w)\ge -C|t|^{q+1}
\end{equation}
for some curve-dependent $C$ near zero.

Suppose the first nonzero derivative of $F\circ\gamma$ through order $q$ has
order $k$. Its Taylor expansion begins with $bt^k$ for $b\ne0$. If $k$ is
odd, one of the two signs of $t$ makes $bt^k<0$ at order $k$. If $k$ is even
and $b<0$, both signs do so. Either case contradicts
\eqref{eq:lower-order-bound}, since $k\le q$. The first nonzero derivative is
therefore even and positive. If none exists, all derivatives through $q$
vanish. This is $q$-jet criticality.
\end{proof}

\begin{proof}[Proof of Theorem~\ref{thm:threshold}]
First-order criticality follows from $q$-jet criticality, so
Lemma~\ref{lem:first-order} supplies active-gradient equality. If an inactive
$j$ violated $g_j\ge\lambda$, then $r_j\le q$ and
Lemma~\ref{lem:boundary-jet} would give a negative first nonzero derivative of
even order $r_j$ along $\gamma_j$. This contradicts $q$-jet criticality.
Thus all simplex KKT inequalities hold.

For sharpness, fix an inactive $j$ with $r_j>q$ and take $f(x)=-x_j$. For
every $C^q$ curve $\gamma$ through $\bar w$,
Lemma~\ref{lem:inactive-order} gives
\[
  (f\circ\Phi)(\gamma(t))=-\Phi_j(\gamma(t))=O(|t|^{r_j})=o(|t|^q).
\]
All derivatives through order $q$ vanish, so $\bar w$ is $q$-jet critical.
At $\bar x$, however, $g_j=-1$ and $g_i=0$ on $S$, which violates the
missing KKT inequality. In the full-support case, Lemma~\ref{lem:first-order}
already supplies every KKT equality and no inequality is missing.
\end{proof}

\section{A boundary-jet certificate}\label{sec:certificate}

Identity \eqref{eq:jet-identity} is useful independently of the threshold
theorem. It gives a direct formula for a reduced cost from a single boundary
jet. If the contact tensor is known, then
\begin{equation}\label{eq:reduced-cost}
  g_j-\lambda
  =\frac{A_0}{D^{r_j}a_j(\bar w_j)[v_j^{r_j}]}
   (F\circ\gamma_j)^{(r_j)}(0).
\end{equation}
Thus the sign of the missing simplex KKT inequality is encoded in the first
nonzero boundary derivative. This certificate does not require convexity.

The normalization in \eqref{eq:lift} is essential to the coefficient. Adding
raw mass to coordinate $j$ simultaneously scales down all active coordinates.
The combination $g_j-g^\top\bar x$, rather than $g_j$ alone, is the resulting
directional derivative. Once the active derivatives agree,
$g^\top\bar x=\lambda$ and the coefficient becomes the usual reduced cost.

\section{Deep squared Hadamard products}\label{sec:hadamard}

Consider $L$ vector factors $u_1,\ldots,u_L\in\R^n$ and define
\begin{equation}\label{eq:deep-lift}
  a_i(u)=\prod_{\ell=1}^L u_{\ell i}^2,
  \qquad
  \Phi_i(u)=\frac{\prod_{\ell=1}^L u_{\ell i}^2}
  {\sum_{k=1}^n\prod_{\ell=1}^L u_{\ell k}^2}.
\end{equation}
This is block-separable after grouping the parameters by coordinates as
$w_i=(u_{1i},\ldots,u_{Li})\in\R^L$.
For a fixed coordinate $i$, let
\[
  z_i=\#\{\ell:u_{\ell i}=0\}
\]
be its zero multiplicity.

\begin{proposition}[Zero multiplicity equals half the contact order]
\label{prop:zero-multiplicity}
At an inactive coordinate $i$ for \eqref{eq:deep-lift}, the contact order is
\[
  r_i=2z_i.
\]
At an active coordinate, the block gradient of $a_i$ is nonzero.
\end{proposition}

\begin{proof}
For every zero factor, write its perturbation as $u_{\ell i}(t)=t v_{\ell i}$.
Keep every nonzero factor fixed. If all chosen $v_{\ell i}$ are nonzero, then
\[
  a_i(u(t))
  =t^{2z_i}
   \left(\prod_{\ell:u_{\ell i}=0}v_{\ell i}^2\right)
   \left(\prod_{\ell:u_{\ell i}\ne0}u_{\ell i}^2\right).
\]
The coefficient is positive. Every monomial in $a_i$ contains two copies of
each zero factor, so no derivative of total order below $2z_i$ can be
nonzero. This proves the contact order. If $a_i(u)>0$, every factor is
nonzero and
\[
  \frac{\partial a_i}{\partial u_{\ell i}}
  =2u_{\ell i}\prod_{p\ne\ell}u_{p i}^2\ne0,
\]
which proves active nondegeneracy.
\end{proof}

\begin{corollary}[Which missing coordinates the Hessian sees]
\label{cor:hessian-visibility}
Let $f$ be convex and twice continuously differentiable. A point of the deep
lift \eqref{eq:deep-lift} is an ordinary second-order critical point of
$f\circ\Phi$ if and only if
\begin{align*}
  g_i&=\lambda &&(i\in S),\\
  g_j&\ge\lambda &&(j\notin S\text{ and }z_j=1).
\end{align*}
No KKT inequality is tested by the Hessian when $z_j\ge2$.
More generally, $q$-jet criticality tests coordinate $j$ exactly when
$2z_j\le q$.
\end{corollary}

\begin{proof}
Combine Proposition~\ref{prop:zero-multiplicity} with
Theorem~\ref{thm:convex}.
\end{proof}

The corollary distinguishes two boundary strata that look identical in
simplex space. A missing coordinate with one zero factor is visible at second
order. The same missing coordinate with several zero factors is invisible to
the Hessian because its raw mass has flatter contact with the boundary.

\section{Even-power lifts and exponentially many false second-order points}
\label{sec:powers}

The tied even-power lift
\begin{equation}\label{eq:power-lift}
  \Phi_i(u)=\frac{u_i^{2m}}{\sum_k u_k^{2m}},
  \qquad m\ge1,
\end{equation}
has contact order $2m$ at every zero coordinate. When $m=1$, ordinary
second-order stationarity sees every missing KKT inequality. When $m\ge2$,
the Hessian sees none of them.

The loss is substantial even for a strongly convex objective. Let
\begin{equation}\label{eq:quadratic}
  f(x)=\frac12\left\lVert x-\frac1n\mathbf1\right\rVert^2.
\end{equation}
Its unique simplex minimizer is the uniform point
$x^\star=\mathbf1/n$. For every nonempty proper subset
$S\subsetneq\{1,\ldots,n\}$, define the face-uniform point
\[
  x^S_i=
  \begin{cases}
    1/|S|,&i\in S,\\
    0,&i\notin S.
  \end{cases}
\]
and choose any $u^S$ with $u_i^{2m}$ constant and positive on $S$ and zero
off $S$.

\begin{proposition}[Exponential false second-order family]
\label{prop:exponential}
For every $m\ge2$, each $u^S$ is an ordinary second-order critical point of
$f\circ\Phi$ under \eqref{eq:power-lift}, while $x^S$ is not simplex
stationary. Hence the lift has at least $2^n-2$ distinct nonglobal
second-order critical images.
\end{proposition}

\begin{proof}
For \eqref{eq:quadratic},
\[
  g_i(x^S)=x^S_i-1/n.
\]
This quantity is constant on $S$, so active-gradient equality holds. Every
inactive coordinate has contact order $2m>2$. Theorem~\ref{thm:convex} with
$q=2$ therefore makes $u^S$ second-order critical.

The active value is $\lambda=1/|S|-1/n>0$, whereas every inactive derivative
is $g_j=-1/n<\lambda$. Thus every missing KKT inequality fails and $x^S$ is
not stationary. Distinct nonempty proper subsets have distinct supports, so
their images are distinct. There are $2^n-2$ such subsets.
\end{proof}

Proposition~\ref{prop:exponential} concerns critical points rather than local
minima. Along a curve that activates a missing coordinate, the objective
decreases first at order $2m$. The example shows why a Hessian-level landscape
statement cannot be transferred uniformly through a flatter boundary lift.

\section{Closest known results and novelty boundary}\label{sec:related}

Li, McKenzie, and Yin analyze the Hadamard square parametrization of simplex
and related constraints, including the correspondence between second-order
stationarity in parametrized space and first-order stationarity in the
original problem \cite{LiMcKenzieYin2023}. Levin, Kileel, and Boumal give a
general theory of first- and second-order landscape preservation under smooth
parametrizations \cite{LevinKileelBoumal2024}. Those results cover the standard
square lift and supply the immediate order-two precedent for this work.

Tang and Toh study Hadamard parametrization for composite optimization, while
Ouyang and coauthors use Hadamard difference parametrizations in sparse
optimization \cite{TangToh2024,OuyangEtAl2024}. Achour, Malgouyres, and
Gerchinovitz classify second-order critical behavior in deep linear networks,
another setting in which factor multiplicity changes the critical landscape
\cite{AchourMalgouyresGerchinovitz2021}.

Kolb, M\"uller, Bischl, and R\"ugamer give a 2026 synthesis of Hadamard
overparametrization for sparse regularization, including deeper product maps
and preservation of local minima through openness
\cite{KolbEtAl2026}. Their framework makes depth an explicit modeling choice.
It does not provide the finite-jet simplex KKT filtration proved here.

Levin, Kileel, and Boumal explicitly list the general cost-independent
``$k\Rightarrow1$'' problem as a future direction
\cite{LevinKileelBoumal2024}. The precise claim advanced here is narrower than
a general landscape theory.
For normalized block-separable nonnegative simplex lifts, contact order gives
an exact convex characterization at every finite jet, a sharp universal jet
threshold for preservation of first-order simplex stationarity, and the
zero-multiplicity rule $r_i=2z_i$ for deep squared Hadamard products. The
standard square case $r_i=2$ is prior art and is included only as the first
member of the order filtration.

This novelty claim remains subject to independent literature review. In
particular, higher-order optimality conditions, singular reparametrizations,
and algebraic treatments of Hadamard-product models use several different
terminologies. A result equivalent to the contact-order filtration may exist
outside the references identified here.

\section{Limitations and open questions}\label{sec:limitations}

The analysis has several deliberate limits.
\begin{enumerate}[label=(\roman*)]
  \item Every inactive block is assumed to have finite contact order. Flat
  smooth functions with all derivatives zero at the boundary are invisible to
  every finite jet and require a non-Taylor analysis.
  \item Active nondegeneracy is essential. If an active raw weight is itself
  critical, first-order lifted stationarity need not force equality of active
  objective derivatives.
  \item Convexity is used for the sufficiency half of the exact finite-jet
  characterization. The sharp universal threshold recovers first-order KKT
  stationarity for nonconvex objectives but does not characterize all
  nonconvex higher-order terms.
  \item The theorem is local and qualitative. It does not give iteration
  complexity, conditioning bounds, or an algorithm for escaping a high-order
  flat saddle.
  \item The framework addresses normalized block-separable lifts. Coupled raw
  weights may mix inactive coordinates in their first nonzero homogeneous
  terms and require a cone-valued contact invariant.
\end{enumerate}

Natural next questions include a coordinate-free extension to coupled maps,
quantitative error bounds based on the smallest positive contact tensor, and
algorithmic mechanisms that detect reduced costs without differentiating to
order $r_{\max}$.

\section{Conclusion}

The singularity of a simplex lift at the boundary has an order, and that order
controls which constrained stationarity conditions are visible. For
normalized block-separable nonnegative lifts, the rule is exact: a finite jet
tests precisely the inactive KKT inequalities whose contact orders it reaches.
Deep Hadamard products turn zero multiplicity into contact order, making the
gap between shallow and deep parametrizations explicit. This perspective
provides a local diagnostic for landscape-transfer claims and a design rule
for simplex parametrizations: boundary flatness beyond order two creates
stationary behavior that the Hessian cannot detect.

\bibliographystyle{plain}
\bibliography{paper/references}

\end{document}
